\documentclass{article}\usepackage[T1]{fontenc}\usepackage{lmodern}
\usepackage[margin=16mm]{geometry}
\usepackage{gradbars}
\setlength{\parindent}{0pt}
\usegradbarslibrary{reports,legend}
\gradscale{review-range}{min=0,max=100}
\graddeclaremetric{review}{title={Review score},role=score,
 unit={\,pt},better=higher,precision=1,
 scale group=review-range,aggregate=mean}
\gradbarsloadcsv{reviews}{docs/data/report-metrics.csv}
\gradtableschema{review-columns}{
 \gradcolumn[source=Project,width=28mm]{project}{Project}
 \gradmetricbind[source=Review,mark=both]{review}{score}
 \gradmetriccardcolumn[source=History,
  options={height=5mm,label width=10mm}]{review}{history}}
\begin{document}
\section*{Report metrics across tables}
Fictitious data; one metric freezes the range, unit and direction.\par
\medskip
\gradmetricuse[width=40mm]{review}{80}\quad
\gradmetricnumber{review}{0}\par\medskip
\gradmetriccard{review}{0,40,NA,80}\par\medskip
\gradmetriccard{review}{20,40,NA}\par\medskip
\textbf{Table 1: top two; visible mean 60.0 pt}\par
\begin{gradtable}[data=reviews,style=metric-report,width=\linewidth,
 sort=score,order=descending,limit=2]
 \graduseschema{review-columns}\gradsummary{Mean}
\end{gradtable}\par\medskip
\textbf{Table 2: follow-up values; data mean 50.0 pt}\par
\begin{gradtable}[data=reviews,style=metric-report,width=\linewidth,
 limit=2]
 \gradcolumn[source=Project,width=28mm]{project}{Project}
 \gradmetricbind[source=Followup,mark=both]{review}{score}
 \gradsummary[scope=data]{Mean}
\end{gradtable}\par
The current value is the last history position; final NA does not fall back. Table 2 shows two rows but summarizes all valid follow-up values.
\end{document}
